% !TeX root = ../../Book2.tex
\section{非齐次方程}
\label{b2:sec:0902}

外力在每一时刻都改变速度。线性叠加允许把持续作用拆成各个时刻产生的初速度增量，再分别传播到观察时刻。积分上限处的贡献决定了非齐次项；若只对积分内部求导，便会漏掉外力本身。

\subsection{Duhamel 原理}
\label{b2:subsec:090201}

取 \(c>0\)、\(h\in C(\mathbb R;\mathbb C)\)、\(t\in\mathbb R\)。先把零位移、给定速度的解算子记为
\[
 S_c(t)h(x)\coloneqq\dfrac1{2c}\int_{x-ct}^{x+ct}h(y)\,\mathrm{d}y.
\]
负时间按有向积分解释。它满足 \(S_c(0)h=0\)、\(\partial_tS_c(0)h=h\)。若时刻 \(s\) 的力为 \(F(s,\cdot)\)，它在时长 \(t-s\) 内产生的响应是 \(S_c(t-s)F(s,\cdot)\)。对这些响应积分的做法称为\term[duhamel-yuanli]{Duhamel 原理}[Duhamel's principle]。
\symidx{\(S_c(t)\)：一维波方程的初速度传播算子}

\begin{theorem}[Duhamel，一维波方程]\label{b2:thm:09-duhamel}
设 \(c>0\)、\(T>0\)，\(f\in C^2(\mathbb R;\mathbb C)\)、\(g\in C^1(\mathbb R;\mathbb C)\)、\(F\in C^1([0,T]\times\mathbb R;\mathbb C)\)。问题
\[
 u_{tt}-c^2u_{xx}=F,\quad u(0)=f,\quad u_t(0)=g
\]
在 \([0,T]\times\mathbb R\) 上有唯一 \(C^2\) 经典解；它等于\eqref{b2:eq:09-dalembert}的齐次解加上
\begin{equation}\label{b2:eq:09-duhamel}
 w(t,x)=\int_0^tS_c(t-s)F(s,\cdot)(x)\,\mathrm{d}s
 =\dfrac1{2c}\int_0^t\int_{x-c(t-s)}^{x+c(t-s)}F(s,y)\,\mathrm{d}y\,\mathrm{d}s.
\end{equation}
\end{theorem}
\begin{proof}
记 \(V(\tau,s,x)=S_c(\tau)F(s,\cdot)(x)\)。固定 \(s\) 时，\(V_{\tau\tau}=c^2V_{xx}\)，且 \(V(0,s,x)=0\)、\(V_\tau(0,s,x)=F(s,x)\)。各阶求导只涉及一个紧三角形上的 \(F\)、\(F_x\)，故可在积分号下求导。Leibniz 公式给出
\[
 w_t=\int_0^tV_\tau(t-s,s,x)\,\mathrm{d}s,\quad
 w_{tt}=F(t,x)+\int_0^tV_{\tau\tau}(t-s,s,x)\,\mathrm{d}s.
\]
另一方面，\(c^2w_{xx}\) 正是第二式的积分项。所以 \(w_{tt}-c^2w_{xx}=F\)。积分区间在 \(t=0\) 退化，给出 \(w(0)=w_t(0)=0\)。减去这个特解后使用\cref{b2:thm:09-dalembert}，即可得到一般初值的存在性与唯一性。
\end{proof}

外力项的积分区域是观察点向过去展开的三角形。对固定 \(s\)，水平截面宽度为 \(2c(t-s)\)：较早施加的力能够从较远的位置传来，刚施加的力只能从观察点附近传来。这一几何结构比公式中的常数更有普遍性，后面的高维问题将把区间改成球。

\Needspace{6\baselineskip}
\begin{proposition}[初速度传播的平方范数界]\label{b2:prop:09-velocity-propagator}
设 \(c>0\)，\(t\in\mathbb R\)，\(h\in L^2(\mathbb R;\mathbb C)\)。记 \(H^1(\mathbb R;\mathbb C)=\{v\in L^2(\mathbb R;\mathbb C):v'\in L^2(\mathbb R;\mathbb C)\}\)，其中导数为弱导数。以有向积分定义
\[
 S_c(t)h(x)=\frac1{2c}\int_{x-ct}^{x+ct}h(y)\,\mathrm dy.
\]
则 \(S_c(t)h\in H^1(\mathbb R)\)，并有
\[
 \|S_c(t)h\|_2\leq|t|\|h\|_2,
 \quad c\|\partial_xS_c(t)h\|_2\leq\|h\|_2.
\]
映射 \(t\mapsto S_c(t)h\) 属于 \(C(\mathbb R;H^1)\cap C^1(\mathbb R;L^2)\)，且
\[
 \partial_tS_c(t)h=\frac{h(\cdot+ct)+h(\cdot-ct)}2,
 \quad\partial_xS_c(t)h=\frac{h(\cdot+ct)-h(\cdot-ct)}{2c}.
\]
\end{proposition}
\begin{proof}
在有限区间上，\(L^2\) 函数属于 \(L^1\)，故区间积分不依赖几乎处处代表。对 \(t>0\)，它是与 \((2c)^{-1}\mathbf1_{[-ct,ct]}\) 的卷积，此核的 \(L^1\) 范数为 \(t\)，Young 不等式给出零阶界；负时间由有向积分给出相同的 \(|t|\) 界。

先对光滑函数作端点求导，得到所述两式。以 \(C_c^\infty\) 函数在 \(L^2\) 中逼近 \(h\)，平移的等距性使两式右侧在 \(L^2\) 中收敛。空间分部积分取极限，识别第二式为弱导数，给出其范数界。对于时间导数，也可直接在 \(L^2\) 中写成
\[
 S_c(t)h=\frac12\int_0^t
 \bigl(h(\cdot+cs)+h(\cdot-cs)\bigr)\,\mathrm{d}s.
\]
平移强连续性使被积函数在 \(L^2\) 中连续，向量值微积分基本定理给出第一式。空间导数公式的平移连续性又给出 \(H^1\) 连续性，包含 \(t=0\) 时空间导数为零的结论。
\end{proof}

对 \(c>0\)、\(f\in L^2(\mathbb R;\mathbb C)\)，定义初位移算子
\[
 C_c(t)f=\frac{f(\cdot+ct)+f(\cdot-ct)}2,\quad t\in\mathbb R.
\]
平移的等距性给出 \(\|C_c(t)f\|_2\leq\|f\|_2\)，平移强连续性给出 \(t\mapsto C_c(t)f\) 的 \(L^2\) 连续性。若 \(f\in H^1(\mathbb R;\mathbb C)\)，平移与弱导数交换，故还有
\[
 \partial_xC_c(t)f=C_c(t)f',\quad
 \|C_c(t)f\|_{H^1}\leq\|f\|_{H^1},\quad
 \partial_tC_c(t)f=\frac c2\bigl(f'(\cdot+ct)-f'(\cdot-ct)\bigr).
\]
这里 \(H^1\) 范数平方为 \(\|f\|_2^2+\|f'\|_2^2\)。第一式可将平移后的函数与光滑紧支集测试函数配对，再换元和分部积分验证；它与平移连续性一起给出 \(H^1\) 连续性。最后一式是在 \(L^2\) 中的强导数：\cref{b2:prop:04-difference-quotients}说明，具有平方可积弱导数的函数，其差商在 \(L^2\) 中趋于该弱导数，且范数由 \(\|f'\|_2\) 控制。将差商分别用于两个平移即得所示时间导数；不能对一般等价类直接作逐点求导。

\ProseParagraphBreak

对粗外力，只要求时间上的可积性。记 \(F\in L^1((0,T);L^2(\mathbb R;\mathbb C))\) 是指 \(F\) 为强可测的 \(L^2\) 值函数，且 \(\int_0^T\|F(s)\|_2\,\mathrm{d}s<\infty\)。这里强可测意为几乎处处是有限值简单函数列的范数极限；向量值积分由简单函数在积分范数中的逼近定义，积分三角不等式使它不依赖逼近的选择。传播算子恰好把这一外力空间送入后面能量方法所需的解空间。

\Needspace{6\baselineskip}
\begin{corollary}[表示式的正则性与数据连续依赖]\label{b2:cor:09-propagator-regularity}
设 \(c>0\)、\(T>0\)，\(f\in H^1(\mathbb R;\mathbb C)\)、\(g\in L^2(\mathbb R;\mathbb C)\)，且 \(F\in L^1((0,T);L^2(\mathbb R;\mathbb C))\)。定义
\[
 u(t)=C_c(t)f+S_c(t)g+w(t),\quad
 w(t)=\int_0^tS_c(t-s)F(s)\,\mathrm{d}s,
\]
其中 \(C_c(t)f=(f(\cdot+ct)+f(\cdot-ct))/2\)，\(S_c(t)h=(2c)^{-1}\int_{\cdot-ct}^{\cdot+ct}h(y)\,\mathrm{d}y\)，外力积分在 \(H^1\) 中取值。则
\[
 u\in C([0,T];H^1(\mathbb R;\mathbb C))\cap C^1([0,T];L^2(\mathbb R;\mathbb C)),
\]
且其强时间导数为
\[
 u_t(t)=\frac c2\bigl(f'(\cdot+ct)-f'(\cdot-ct)\bigr)
       +C_c(t)g+\int_0^tC_c(t-s)F(s)\,\mathrm{d}s.
\]
特别，\(u(0)=f\)、\(u_t(0)=g\)。对 \(0\leq t\leq T\)，有
\[
 \begin{aligned}
 \|u(t)\|_2&\leq\|f\|_2+t\|g\|_2
               +\int_0^t(t-s)\|F(s)\|_2\,\mathrm{d}s,\\
 \|u_x(t)\|_2&\leq\|f'\|_2+c^{-1}\|g\|_2
               +c^{-1}\int_0^t\|F(s)\|_2\,\mathrm{d}s,\\
 \|u_t(t)\|_2&\leq c\|f'\|_2+\|g\|_2
               +\int_0^t\|F(s)\|_2\,\mathrm{d}s.
 \end{aligned}
\]
因此表示式给出从 \(H^1\times L^2\times L^1((0,T);L^2)\) 到所列解空间的连续线性映射；解空间采用 \(\sup_{0\leq t\leq T}(\|u(t)\|_{H^1}+\|u_t(t)\|_2)\) 范数。
\end{corollary}
\begin{proof}
初位移与初速度两项的结论已由上述平移计算和\cref{b2:prop:09-velocity-propagator}得到。只需证明外力积分的正则性。对 \(0\leq s\leq t\leq T\)，
\[
 \|S_c(t-s)F(s)\|_{H^1}
 \leq\bigl(T+c^{-1}\bigr)\|F(s)\|_2.
\]
用 \(L^2\) 值简单函数逼近 \(F\)，再利用 \(S_c\) 的强连续性，得到被积函数的 \(H^1\) 强可测性；所示可积控制保证积分存在。这里 \(H^1\) 完备：若函数和弱导数分别在 \(L^2\) 中收敛，对测试函数分部积分取极限，便识别出极限的弱导数。因而简单函数积分的 Cauchy 极限仍在 \(H^1\) 中。

比较两个相邻时刻，在公共积分区间上使用 \(S_c\) 的 \(H^1\) 强连续性和上述可积控制，在新增或删去的短区间上使用 \(\int\|F(s)\|_2\) 的绝对连续性，得到 \(w\in C([0,T];H^1)\)。时间导数为
\[
 w_t(t)=\frac12\int_0^t
 \bigl(F(s,\cdot+c(t-s))+F(s,\cdot-c(t-s))\bigr)\,\mathrm{d}s
 \quad\text{于 }L^2.
\]
具体地，设时间增量为 \(\delta\)。在公共积分区间上，\(S_c(t+\delta-s)-S_c(t-s)\) 除以 \(\delta\) 后作用于 \(F(s)\)，由传播算子的强导数公式趋于 \(C_c(t-s)F(s)\)。差商的 \(L^2\) 范数不超过 \(\|F(s)\|_2\)：把它写为 \(C_c(r)F(s)\) 在长度 \(|\delta|\) 的时间区间上的平均即可。因此可在积分中取极限。在新增或删去的短区间上，\(S_c\) 的时间参数满足 \(|\tau|\leq|\delta|\)，由 \(\|S_c(\tau)F(s)\|_2\leq|\tau|\|F(s)\|_2\)，该部分除以 \(|\delta|\) 后的范数至多为短区间上的 \(\int\|F(s)\|_2\)，趋于零。于是正、负差商给出同一个导数；端点使用单侧差商。

所得导数中，\(C_c(t-s)F(s)\) 对 \(t\) 在 \(L^2\) 中强连续，范数仍由 \(\|F(s)\|_2\) 控制。再次分拆公共区间与短区间，得到 \(w_t\in C([0,T];L^2)\)，故 \(w\in C^1([0,T];L^2)\)。初值由 \(S_c(0)=0\)、\(C_c(0)=I\) 及积分在零时刻消失得到。

最后对三项分别使用传播算子界与积分三角不等式，得到所列三个估计。把同一估计用于两组数据之差，并令 \(t\leq T\)，便得到解及其空间、时间导数在整个时间区间上的一致范数连续依赖。
\end{proof}

这一步不要求 \(F\) 有时间导数；粗数据对应的二阶时间导数将在\secref{b2:sec:0903}的分布方程中解释。上述正则性先为该节的截断能量计算提供适用条件，随后再由能量方法证明解的唯一性与更精确的稳定性估计。

\subsection{外力项}
\label{b2:subsec:090202}

设 \(c>0\)、\(T>0\)、\(a\in C^1([0,T];\mathbb C)\)。若 \(F(s,y)=a(s)\) 与空间无关，且初值为零，公式化为
\[
 u(t,x)=\int_0^t(t-s)a(s)\,\mathrm{d}s.
\]
因此常数外力 \(a(s)\equiv a_0\) 产生 \(u=a_0t^2/2\)，恰与没有空间变化时的常微分方程 \(u_{tt}=a_0\) 一致。常数 \(1/(2c)\) 被区间长度中的 \(2c\) 抵消，这也是检查归一化的一种简便方法。本例属于不要求有限总能量的经典解情形：非零空间常数不属于 \(L^2(\mathbb R)\)，故只要 \(a\not\equiv0\)，连续外力 \(F\) 就不属于 \(L^1((0,T);L^2(\mathbb R))\)；相应位移在取非零常数值的时刻也不属于 \(H^1(\mathbb R)\)。

\ProseParagraphBreak

外力的大小受到控制，并不意味着位移在所有时间上一致有界。取 \(c>0\)、\(k>0\)、\(\omega=ck\)，令
\[
 F(t,x)=\cos(\omega t)\cos(kx),\quad f=g=0.
\]
设 \(u(t,x)=a(t)\cos(kx)\)，则 \(a''+\omega^2a=\cos(\omega t)\)，从而
\[
 a(t)=\dfrac{t\sin(\omega t)}{2\omega}.
\]
直接求导可核对方程与零初值。外力的时间频率与自由振动频率相同，持续作用可使振幅增长；这种现象称为\term[gongzhen]{共振}[resonance]。本例为空间周期的经典解，不具有整条实线上的有限总能量，因此不把它代入稍后的全空间能量积分。

\ProseParagraphBreak

若 \(c>0\)、\(T>0\)，且 \(F\in C^1([0,T]\times\mathbb R;\mathbb C)\) 在该时间带上一致有界，零初值响应 \(w\) 由\eqref{b2:eq:09-duhamel}直接满足
\[
 \|w(t)\|_\infty\leq\int_0^t(t-s)\|F(s)\|_\infty\,\mathrm{d}s.
\]
在 \(F\) 一致有界时右端不超过 \(t^2\|F\|_\infty/2\)。该估计没有利用振荡抵消，适合控制局部存在时间内的误差；振荡外力的长时间响应还取决于相位抵消。

\subsection{有限传播速度}
\label{b2:subsec:090203}

给定传播速度 \(c>0\)、\(t_0>0\)、\(x_0\in\mathbb R\)，记过去的传播区域为
\[
 K_{t_0,x_0}\coloneqq
 \{\,(s,y):0\leq s\leq t_0,\ |y-x_0|\leq c(t_0-s)\,\}.
\]
若称某个区域为观察点的\term[yilai-quyu]{依赖区域}[domain of dependence]，意思是该处的解只由区域内的数据确定。相反，从一个初始位置出发可能受到影响的未来区域称为\term[yingxiang-quyu]{影响区域}[domain of influence]。这两个说法分别固定终点和起点，不能把箭头方向混同。

\begin{proposition}[有限传播速度]\label{b2:prop:09-finite-speed}
设 \(c>0\)、\(T>0\)、\(0\leq t_0\leq T\)、\(x_0\in\mathbb R\)，记 \(K_{t_0,x_0}=\{(t,x):0\leq t\leq t_0,\ |x-x_0|\leq c(t_0-t)\}\)。设开集 \(U\subset\mathbb R^2\) 包含该闭三角形，\(u,\widetilde u\in C^2(U;\mathbb C)\) 分别满足
\[
 u_{tt}-c^2u_{xx}=F,\quad
 \widetilde u_{tt}-c^2\widetilde u_{xx}=\widetilde F,
\]
其中 \(F,\widetilde F\in C(U;\mathbb C)\)。若两解的初始位移、初始速度在 \(I=[x_0-ct_0,x_0+ct_0]\) 上分别相等，且 \(F=\widetilde F\) 于 \(K_{t_0,x_0}\)，则两解在整个 \(K_{t_0,x_0}\) 上相等。

\ProseParagraphBreak

特别，设 \(u\in C^2([0,T]\times\mathbb R;\mathbb C)\) 满足全空间方程 \(u_{tt}-c^2u_{xx}=F\)，\(F\) 连续。若其初始位移、初始速度及每个时刻的外力都支撑在固定闭集 \(A\subset\mathbb R\) 内，则
\[
 \operatorname{supp}u(t,\cdot)\subseteq
 \{\,x:\operatorname{dist}(x,A)\leq ct\,\},\quad 0\leq t\leq T.
\]
\end{proposition}
\begin{proof}
当 \(t_0=0\) 时，结论就是给定的初始位移相等。以下设 \(t_0>0\)，记 \(v=u-\widetilde u\)、\(G=F-\widetilde F\)，并在 \(U\) 上令
\[
 p=v_t+cv_x,\quad q=v_t-cv_x.
\]
由混合偏导可交换，
\[
 p_t-cp_x=G,\quad q_t+cq_x=G.
\]
取任意 \((t,x)\in K_{t_0,x_0}\)。对 \(0\leq s\leq t\)，有
\[
 |x\pm c(t-s)-x_0|
 \leq|x-x_0|+c(t-s)\leq c(t_0-s).
\]
所以从该点向初值线追踪的两条特征线段始终留在原三角形内，无须把局部解延拓到全实线。沿 \(s\mapsto(s,x+c(t-s))\) 和 \(s\mapsto(s,x-c(t-s))\) 分别积分，得到
\[
 \begin{aligned}
 p(t,x)&=p(0,x+ct)+\int_0^tG(s,x+c(t-s))\,\mathrm{d}s,\\
 q(t,x)&=q(0,x-ct)+\int_0^tG(s,x-c(t-s))\,\mathrm{d}s.
 \end{aligned}
\]
底边上 \(v(0,\cdot)=0\)，故其空间导数在 \(I\) 内部为零，并由连续性在两端点也为零；又有 \(v_t(0,\cdot)=0\)，于是 \(p(0,\cdot)=q(0,\cdot)=0\) 于 \(I\)。特征线上的 \(G\) 也为零，因此 \(p=q=0\) 于整个三角形，进而 \(v_t=v_x=0\)。对所取点，竖直线段 \(\{(s,x):0\leq s\leq t\}\) 也留在三角形内，故
\[
 v(t,x)=v(0,x)+\int_0^tv_t(s,x)\,\mathrm{d}s=0.
\]
因此，特征三角形内的解由其底边数据与内部外力决定。

对于支集结论，将全空间解与零解比较。若 \(\operatorname{dist}(x,A)>ct\)，则该点过去三角形的每个截面都不与 \(A\) 相交，故底边初值和三角形内外力均为零，刚证的局部结论给出 \(u(t,x)=0\)。当 \(A\ne\varnothing\) 时，距离函数连续，使这一结论在 \(x\) 的空间邻域也成立；若 \(A=\varnothing\)，数据全为零，同一局部结论直接给出零解，并约定到空集的距离为 \(+\infty\)。于是均得到所陈述的支集包含关系。
\end{proof}

\cref{b2:fig:09-cone}画出了同一个区域的两种用途：计算 \((t_0,x_0)\) 的位移时读取过去的数据；证明唯一性时则比较一整个收缩区间上的能量。第二种方式不依赖显式公式，因而能够推广到其他维数。
\begin{figure}[htbp]
\centering
\begin{tikzpicture}[x=1cm,y=1cm,>=stealth]
\fill[AnalysisBlue!9] (-3,0)--(0,3)--(3,0)--cycle;
\draw[->] (-3.8,0)--(3.8,0) node[right]{\(x\)};
\draw[->] (-3.6,-.2)--(-3.6,3.5) node[above]{\(t\)};
\draw[thick,AnalysisBlue] (-3,0)--(0,3)--(3,0);
\draw[dashed] (-1.8,1.2)--(1.8,1.2);
\fill (0,3) circle(1.7pt) node[above right]{\((t_0,x_0)\)};
\node[below] at (-3,0){\(x_0-ct_0\)};
\node[below] at (3,0){\(x_0+ct_0\)};
\node[below] at (0,0){\(x_0\)};
\node at (0,.6){\(K_{t_0,x_0}\)};
\node[above] at (0,1.2){时刻 \(s\) 的截面};
\end{tikzpicture}
\caption[一维波动方程的依赖三角形]{一维波动方程的依赖三角形。截面半径为 \(c(t_0-s)\)，两侧边界随时间以速度 \(c\) 向观察点收缩；图中横纵比例仅为示意。}
\label{b2:fig:09-cone}
\end{figure}


传播速度的结论只给出支集的外界，不保证界内处处非零。左右行波可以抵消，数据本身也可能具有空隙。此外，在一维即使初始数据紧支撑，固定位置也未必在波前经过以后回到零位移：若 \(f=0\)、\(g\) 紧支撑，则对足够大的 \(t\)，
\[
 u(t,x)=\dfrac1{2c}\int_{\mathbb R}g(y)\,\mathrm{d}y.
\]
这个尾部不违反有限传播，因为它位于允许的传播区域内部。三维方程在相同条件下具有更强的无尾传播性质，须到\chapref{b2:ch:11}重新证明。
